% !TEX root = ../../main.tex

\section{Implementación en GHC}
\label{sec:solucion-implementacion}

El diseño anterior se implementó directamente en \textsf{GHC}, dentro del procesamiento de expresiones del Renamer. Esta ubicación permite reutilizar la representación del compilador, la información de nombres libres y el algoritmo existente de \texttt{ApplicativeDo}, sin introducir una herramienta externa de reescritura ni operar sobre una representación que ya haya perdido la estructura de la \textit{do-notation}.

\subsection{Ubicación en el Pipeline del Compilador}
\label{sec:solucion-pipeline-ghc}

Las opciones de línea de comandos y los pragmas del módulo configuran, entre otros aspectos, las extensiones de lenguaje que utilizarán las etapas posteriores. A partir de esa configuración, el Parser construye el árbol de sintaxis; el Renamer resuelve los identificadores y sus alcances; el Typechecker verifica e incorpora información de tipos; y el Desugarer traduce la representación tipada hacia Core. Para una expresión, la progresión de fases puede resumirse como:

\begin{ReportExpression}{Fases relevantes del AST de una expresión en GHC.}{expr:fases-ast-ghc}
\[
\begin{aligned}
\text{código fuente}
&\xrightarrow{\mathrm{Parser}}
\texttt{HsExpr GhcPs}
\xrightarrow{\mathrm{Renamer}}
\texttt{HsExpr GhcRn}\\
&\xrightarrow{\mathrm{Typechecker}}
\texttt{HsExpr GhcTc}
\xrightarrow{\mathrm{Desugarer}}
\texttt{Core}.
\end{aligned}
\]
\end{ReportExpression}

El AST de \textsf{GHC} está parametrizado por la fase de compilación. Los parámetros \texttt{GhcPs}, \texttt{GhcRn} y \texttt{GhcTc} distinguen, respectivamente, los nodos parseados, renombrados y verificados por tipos. Una expresión de \textit{do-notation} se representa mediante el constructor \texttt{HsDo}, que contiene el tipo de bloque y su lista de statements. El Renamer especializa el procesamiento según el constructor de cada expresión y transforma un nodo \texttt{HsDo} parseado en su versión renombrada mediante el caso resumido en el Código~\ref{lst:rnexpr-hsdo}.

\begin{lstlisting}[style=haskell,caption={Flujo de renombrado de una expresión \texttt{HsDo} en \texttt{GHC.Rename.Expr}.},label={lst:rnexpr-hsdo}]
rnExpr (HsDo _ do_or_lc (L l stmts))
 = do { ((stmts1, _), fvs1) <-
          rnStmtsWithFreeVars (HsDoStmt do_or_lc) rnExpr stmts
            (\ _ -> return ((), emptyFVs))
      ; (pp_stmts, fvs2) <-
          postProcessStmtsForApplicativeDo do_or_lc stmts1
      ; return (HsDo noExtField do_or_lc (L l pp_stmts),
                fvs1 `plusFV` fvs2) }
\end{lstlisting}

La función \texttt{rnStmtsWithFreeVars} retorna cada statement renombrado junto con las variables libres de su cuerpo. \path{postProcessStmtsForApplicativeDo} comprueba luego que \texttt{ApplicativeDo} esté habilitado, que el contexto corresponda a \texttt{DoExpr} y que el código no se encuentre dentro de un bracket de Template Haskell. \path{rearrangeForApplicativeDo} se invoca solo después de estas comprobaciones. La extensión se incorporó en esta última función, antes de que \texttt{stmtTreeToStmts} reconstruya la lista con nodos \texttt{ApplicativeStmt}.

La elección del Renamer permite, por tanto, trabajar sobre statements que ya poseen identidades únicas y todavía conservan su estructura fuente. Un análisis posterior sobre Core habría requerido reconstruir información de la \textit{do-notation}; uno anterior al renombrado no dispondría de las identidades necesarias para distinguir correctamente binders con la misma escritura textual.

\subsection{Activación Declarativa mediante QualifiedDo}
\label{sec:solucion-qualified-do}

\texttt{QualifiedDo} permite escribir un bloque como \texttt{M.do} y resolver desde el módulo calificador \texttt{M} las operaciones utilizadas por la notación, en vez de usar necesariamente las operaciones habituales de \texttt{Prelude} \cite{GHCQualifiedDo}. La implementación aprovecha esta sintaxis existente para realizar un \textit{opt-in} visible en cada bloque. El Código~\ref{lst:commutative-qualified-do} muestra su forma general.

\begin{lstlisting}[style=haskell,caption={Bloque activado mediante \texttt{Control.Monad.CommutativeDo}.},label={lst:commutative-qualified-do}]
{-# LANGUAGE ApplicativeDo #-}
{-# LANGUAGE QualifiedDo #-}

import qualified Control.Monad.CommutativeDo as CD

pair :: CD.CommutativeMonad m => m a -> m b -> m (a, b)
pair ma mb = CD.do
  a <- ma
  b <- mb
  CD.return (a, b)
\end{lstlisting}

Para soportar esta notación se agregaron dos módulos. \path{GHC.Internal.Control.Monad.CommutativeDo}, dentro de \texttt{ghc-internal}, define la clase marcador \texttt{CommutativeMonad}, el valor \texttt{\_\_commutative\_do\_\_} y las operaciones básicas de \texttt{QualifiedDo}, como \verb|(>>=)|, \verb|(>>)| y \texttt{fail}. También exporta \texttt{return}, \texttt{pure}, \texttt{fmap}, \verb|(<*>)| y \texttt{join}, necesarias cuando el procesamiento de \texttt{ApplicativeDo} transforma el bloque. El módulo público \path{Control.Monad.CommutativeDo}, dentro de \texttt{base}, reexporta esa interfaz para los programas de usuario.

La clase vacía \texttt{CommutativeMonad} representa una afirmación del programador, no una demostración realizada por \textsf{GHC}. Las funciones exportadas por el módulo \texttt{CD} incorporan esa restricción de clase y su uso será verificado durante el typechecking. De manera independiente, durante el renombrado \texttt{isCommutativeQualifiedDo} consulta únicamente si el módulo que califica el bloque exporta un nombre \texttt{\_\_commutative\_do\_\_}. La presencia del marcador habilita la generación de órdenes alternativos; su ausencia conserva el orden original como único candidato. Por ello, el mecanismo completo sigue siendo una declaración de confianza y no una prueba interna de conmutatividad.

La activación efectiva requiere además que \texttt{ApplicativeDo} esté habilitado y que el bloque alcance el procesamiento descrito en la subsección anterior. Aunque el detector contempla también la forma interna de \texttt{MDoExpr}, el gate actual de \texttt{postProcessStmtsForApplicativeDo} aplica la transformación solo sobre \texttt{DoExpr}; por ello, los experimentos utilizan \texttt{CD.do}. La flag de selección presentada al final de esta sección no participa en la activación.

\subsection{Representación de Statements y Construcción del Grafo}
\label{sec:solucion-implementacion-grafo}

La información requerida por el grafo se reúne en \texttt{StmtDepInfo}. Como muestra el Código~\ref{lst:stmt-dep-info}, cada valor conserva el índice estable del statement, el nodo renombrado, sus variables libres originales y los conjuntos locales de lectura y escritura.

\begin{lstlisting}[style=haskell,caption={Información de dependencias almacenada para cada statement.},label={lst:stmt-dep-info}]
data StmtDepInfo = StmtDepInfo
  { sdiIndex       :: !Int
  , sdiStmt        :: !(ExprLStmt GhcRn)
  , sdiFvsOriginal :: !FreeVars
  , sdiReadsLocal  :: !FreeVars
  , sdiWritesLocal :: !FreeVars
  }

type StmtDepNode = DG.Node Int StmtDepInfo
type StmtDepGraph = DG.Graph StmtDepNode
\end{lstlisting}

El tipo \texttt{ExprLStmt GhcRn} corresponde a un statement de expresión localizado y renombrado. \texttt{sdiFvsOriginal} conserva todos los nombres libres entregados por el Renamer, incluidas referencias locales y nombres globales como funciones o constructores. \texttt{sdiReadsLocal} restringe ese conjunto a nombres definidos en el mismo bloque, mientras que \texttt{sdiWritesLocal} contiene los binders introducidos por el statement. Ambos campos usan \texttt{FreeVars}, representación basada en conjuntos de \texttt{Name}s.

El Código~\ref{lst:compute-stmts-info} implementa las definiciones de la Expresión~\ref{expr:lecturas-escrituras-locales}. La utilidad existente \texttt{collectStmtBinders} recorre las distintas formas de statement y, con \texttt{CollNoDictBinders}, recolecta únicamente binders de términos.

\noindent\begin{minipage}{\textwidth}
\begin{lstlisting}[style=haskell,caption={Cálculo de lecturas y escrituras locales en el Renamer.},label={lst:compute-stmts-info}]
computeStmtsInfo stmts = zipWith mkStmtInfo [1 ..] stmts
  where
    all_writes =
      mkNameSet
        (concatMap
          (collectStmtBinders CollNoDictBinders . unLoc . fst)
          stmts)

    mkStmtInfo i (stmt, fvs) = StmtDepInfo
      { sdiIndex       = i
      , sdiStmt        = stmt
      , sdiFvsOriginal = fvs
      , sdiReadsLocal  = fvs `intersectNameSet` all_writes
      , sdiWritesLocal =
          mkNameSet
            (collectStmtBinders CollNoDictBinders (unLoc stmt))
      }
\end{lstlisting}
\end{minipage}

La detección RAW se reduce a una intersección de conjuntos. \texttt{hasDependency} convierte el conjunto resultante en el predicado usado para insertar aristas:

\begin{lstlisting}[style=haskell,caption={Detección de una dependencia RAW entre dos statements.},label={lst:dep-raw}]
depRAW :: StmtDepInfo -> StmtDepInfo -> FreeVars
depRAW si sj =
  sdiWritesLocal si `intersectNameSet` sdiReadsLocal sj

hasDependency :: StmtDepInfo -> StmtDepInfo -> Bool
hasDependency si sj = not (isEmptyNameSet (depRAW si sj))
\end{lstlisting}

En vez de implementar desde cero una estructura de grafo, la solución utiliza \path{GHC.Data.Graph.Directed}. Cada \texttt{DG.DigraphNode} almacena un \texttt{StmtDepInfo}, usa el índice como clave y registra como destinos los statements posteriores que dependen de uno de sus binders. La comprensión del Código~\ref{lst:build-stmt-dep-nodes} corresponde directamente al doble recorrido del Código~\ref{lst:algoritmo-construccion-grafo}.

\noindent\begin{minipage}{\textwidth}
\begin{lstlisting}[style=haskell,caption={Construcción de nodos y aristas mediante \texttt{GHC.Data.Graph.Directed}.},label={lst:build-stmt-dep-nodes}]
buildStmtDepNodes infos = map mkNode infos
  where
    mkNode si = DG.DigraphNode
      { DG.node_payload = si
      , DG.node_key = sdiIndex si
      , DG.node_dependencies =
          [ sdiIndex sj
          | sj <- infos
          , sdiIndex si < sdiIndex sj
          , hasDependency si sj
          ]
      }

buildStmtsDependencyGraph stmts =
  DG.graphFromEdgedVerticesOrd
    $ buildStmtDepNodes
    $ computeStmtsInfo stmts
\end{lstlisting}
\end{minipage}

\subsection{Generación de Permutaciones}
\label{sec:solucion-implementacion-permutaciones}

La función \path{enumerateSemanticTopSortsBounded} recibe el \texttt{StmtDepGraph} y materializa los órdenes descritos en el Código~\ref{lst:algoritmo-enumeracion-topologica}. Para ello construye un \texttt{IntMap} con los grados de entrada, obtiene los sucesores desde el grafo original y los predecesores desde su transposición. En cada llamada recursiva recorre las claves aún no emitidas, conserva aquellas cuyo grado es cero y crea una rama por elección.

La implementación no elimina físicamente los nodos. Al emitir un índice, decrementa en una copia del mapa el grado de entrada de sus sucesores y acumula el nodo seleccionado. Al completar la cantidad original de vértices, invierte el acumulador y extrae el \texttt{StmtDepInfo} almacenado en cada nodo. El recorrido estable de las claves hace que la primera rama reproduzca el orden original, que aparece como candidato cero.

Cada resultado se transforma nuevamente en la secuencia de statements \texttt{ExprLStmt GhcRn} anotados con sus \texttt{FreeVars}, que es la entrada esperada por \texttt{ApplicativeDo}. Las variables libres originales se conservaron precisamente para no recalcular ni perder esta anotación al cambiar el orden de los statements.

\subsection{Integración con ApplicativeDo}
\label{sec:solucion-integracion-applicative-do}

Después de separar el \texttt{LastStmt}, \texttt{rearrangeForApplicativeDo} construye el grafo y enumera las permutaciones solo si \texttt{isCommutativeQualifiedDo} retorna verdadero. En caso contrario crea una lista unitaria con el orden original. La opción \texttt{Opt\_OptimalApplicativeDo} determina luego si se utiliza \texttt{mkStmtTreeHeuristic} o \texttt{mkStmtTreeOptimal}; la función elegida se aplica a cada secuencia.

La representación interna de los planes es el tipo \texttt{StmtTree} del Código~\ref{lst:stmt-tree-cost}. Sus tres constructores corresponden a los tres casos del modelo estructural: un statement elemental, una composición secuencial mediante \textit{bind} y una agrupación de ramas aplicativas. La función \texttt{stmtTreeCost} implementa directamente las reglas de la Expresión~\ref{expr:modelo-costos-estructural}.

\noindent\begin{minipage}{\textwidth}
\begin{lstlisting}[style=haskell,caption={Representación de planes y cálculo de su costo estructural.},label={lst:stmt-tree-cost}]
data StmtTree a
  = StmtTreeOne a
  | StmtTreeBind (StmtTree a) (StmtTree a)
  | StmtTreeApplicative [StmtTree a]

stmtTreeCost :: ExprStmtTree -> Cost
stmtTreeCost (StmtTreeOne _) = 1
stmtTreeCost (StmtTreeBind l r) =
  stmtTreeCost l + stmtTreeCost r
stmtTreeCost (StmtTreeApplicative ts) =
  case ts of
    [] -> 0
    _  -> Partial.maximum (map stmtTreeCost ts)
\end{lstlisting}
\end{minipage}

\texttt{StmtsPermutationInfo} asocia cada árbol con la permutación que lo originó. El índice de candidato parte desde cero, mientras que los índices de \texttt{StmtDepInfo} parten desde uno y conservan las posiciones del programa fuente.

\noindent\begin{minipage}{\textwidth}
\begin{lstlisting}[style=haskell,caption={Información conservada para evaluar y seleccionar candidatos.},label={lst:stmts-permutation-info}]
data StmtsPermutationInfo = StmtsPermutationInfo
  { candIndex :: !Int
  , candPerm  :: [StmtDepInfo]
  , candTree  :: ExprStmtTree
  , candCost  :: !Cost
  }

let stmt_trees = map mkStmtTree all_permutations
let all_permutations_info =
      zipWith3 mkStmtsPermutationInfo
        [0 :: Int ..] validStmtsPerms stmt_trees
let best_candidate = case nonEmpty all_permutations_info of
      Just cs -> minimumBy (comparing candCost) cs
      Nothing -> panic "rearrangeForApplicativeDo: empty stmt_trees"
\end{lstlisting}
\end{minipage}

Conceptualmente, el último bloque aplica un \texttt{map} para evaluar las permutaciones, combina orden y árbol mediante \texttt{zipWith3}, y usa \texttt{minimumBy} para seleccionar el primer mínimo. Un \texttt{filter} separado se utiliza en las trazas para contar cuántos candidatos empatan con el mínimo, pero no interviene en la selección automática.

Una vez elegido el candidato, \texttt{stmtTreeToStmts} reconstruye los statements finales, agrega nuevamente el statement terminal y genera \texttt{ApplicativeStmt} donde el árbol contiene agrupaciones aplicativas. Desde ese punto, el programa continúa por el typechecking y el desazucarado normales de \textsf{GHC}.

\subsection{Selección Forzada para Validación}
\label{sec:solucion-seleccion-forzada}

La política de costo mínimo permite evaluar el resultado optimizado, pero no basta para comprobar individualmente todas las permutaciones producidas. Con ese propósito se agregó la flag experimental

\begin{center}
\texttt{-fado-reorder-candidate-n=<n>}.
\end{center}

\path{GHC.Driver.DynFlags} almacena su valor como \texttt{Maybe Int} y \path{GHC.Driver.Session} registra la opción de línea de comandos. \texttt{rearrangeForApplicativeDo} calcula primero el mejor candidato y luego intenta obtener el elemento indicado por la flag:

\noindent\begin{minipage}{\textwidth}
\begin{lstlisting}[style=haskell,caption={Selección opcional de un candidato por índice.},label={lst:seleccion-forzada-candidato}]
let selected_candidate = do
      i <- reorder_cand_n
      guard (i >= 0)
      listToMaybe (drop i all_permutations_info)

let final_candidate =
      fromMaybe best_candidate selected_candidate
\end{lstlisting}
\end{minipage}

Un índice válido selecciona el candidato base cero solicitado. Si la flag no está presente, el índice es negativo o queda fuera del rango generado, se conserva \texttt{best\_candidate}. Esta opción no habilita la conmutatividad, no crea candidatos adicionales y no altera sus costos; solo reemplaza la decisión final una vez completado el mismo pipeline. El capítulo siguiente utiliza este mecanismo para compilar y ejecutar cada permutación por separado durante la validación semántica.
